\documentclass[11pt]{article}
\usepackage[T1]{fontenc}
\usepackage[a4paper,margin=27mm]{geometry}
\usepackage{lmodern,amsmath,amssymb,amsthm,booktabs,microtype}
\usepackage[round,authoryear]{natbib}
\usepackage{xurl,hyperref}
\hypersetup{hidelinks}
\newtheorem{theorem}{Theorem}[section]
\newtheorem{lemma}[theorem]{Lemma}
\newtheorem{proposition}[theorem]{Proposition}
\newtheorem{corollary}[theorem]{Corollary}
\newcommand{\Q}{\mathbb Q}
\newcommand{\Z}{\mathbb Z}
\newcommand{\C}{\mathbb C}
\newcommand{\F}{\mathbb F}
\newcommand{\Cl}{\operatorname{Cl}}
\newcommand{\Gal}{\operatorname{Gal}}
\newcommand{\Sel}{\operatorname{Sel}}
\newcommand{\N}{\operatorname N}
\newcommand{\workbenchurl}{https://github.com/KokunoYumeto/zeta-and-l-functions-workbench}
\newcommand{\recordurl}{https://github.com/KokunoYumeto/mathematics-workbenches}
\title{Real nonvanishing and strict cyclotomic Selmer groups\\over $\mathbb Q(i)$, after Ishii}
\author{}
\date{24 November 2026}
\begin{document}
\maketitle

\begin{abstract}
Ishii gave explicit theta expressions for the weight-one Artin form of
$\mathbb Q(i,\sqrt[4]{2})$, and Williams and Liu identified this extension
as the conductor-eight ring class field of $\mathbb Q(i)$.
We give an elementary character formula, prove positivity of the theta
function on the imaginary axis by Poisson summation, and deduce strict
positivity of its completed Hecke $L$-function on the entire real axis.
Separately, for every rational prime $p\equiv5\pmod8$, the quartic
subfield has a single prime above $p$ and class number one. A finite
class-field argument proves that its entire cyclotomic $\Z_p$-tower has
trivial $p$-primary class groups. Both faithful quartic-character parts
in the normal closure vanish. For every finite coefficient extension of
$\Q_p$ and every continuous cyclotomic twist, an explicit
finite-coefficient descent proves that the corresponding strict Selmer
group is zero. All finite arithmetic witnesses are printed.
A short comparison proves the same real nonvanishing for the
unramified cubic character of $\mathbb Q(\sqrt{-23})$, using its classical
eta product. An explicit positive Mellin kernel explains why these
positivity arguments alone do not assert RH.
The text and computations were produced by GPT-6 Astra in ChatGPT Pro mode;
no human has checked the text.
\end{abstract}

\section{Context and statements}

Put
\[
 K=\Q(i),\qquad F=\Q(\alpha),\quad \alpha=\sqrt[4]{2}>0,
 \qquad E=K(\alpha).
\]
Write $\Delta=\Gal(E/K)=\langle\delta\rangle$, where
$\delta(\alpha)=i\alpha$, and let $c$ be ordinary complex conjugation.
Thus $\Gal(E/\Q)$ is the dihedral group of order eight,
$c\delta c=\delta^{-1}$, and $E^{\langle c\rangle}=F$.

The analytic object studied below is classical.
\citet[\S2, Remark~1, p.~127]{Ishii1985} gives its separated theta
expression. \citet[Lemma~3.8, p.~332]{WilliamsLiu1994} identifies $E$ with
the ring class field of the order $\Z+8\Z[i]$.
The arguments here fix all character, conductor, and completion
normalizations and give the positivity implication directly.
The class-group argument is the vanishing case of the one-prime method
underlying cyclotomic stability; compare
\citet[Theorem~1(1), p.~264]{Fukuda1994}.
We include its finite proof, so no class-number estimate in a higher
layer is needed.

Let $\chi$ be either faithful complex character of $\Delta$, viewed as
a primitive finite-order Hecke character of $K$ by Artin reciprocity.
The two choices have the same $L$-function. Define
\[
 L(s)=L_K(s,\chi),\qquad
 \Lambda(s)=256^{s/2}(2\pi)^{-s}\Gamma(s)L(s).
 \tag{1}\label{eq:completion}
\]
The omitted factor two in the convention
$\Gamma_{\C}(s)=2(2\pi)^{-s}\Gamma(s)$ is a positive constant.

\begin{theorem}\label{thm:analytic}
The character $\chi$ has conductor $(8)\subset\Z[i]$, of norm $64$.
Its induced two-dimensional Artin representation over $\Q$ has conductor
$256$, determinant $\chi_{-4}$, and
\[
 L(s)=\frac{\zeta_F(s)}{\zeta(s)L(s,\chi_8)}.
 \tag{2}\label{eq:zeta-quotient}
\]
For $q=e^{2\pi i\tau}$, $\Im\tau>0$, its theta series is
\begin{align}
 f(\tau)
 &=\sum_{n\ge1}a_nq^n
 =\frac{\Theta_{1,0,64}(\tau)-\Theta_{4,4,17}(\tau)}2 \notag\\
 &=\left(\frac12\sum_{u\ {\rm odd}}q^{u^2}\right)
       \left(\sum_{v\in\Z}(-1)^vq^{16v^2}\right)
 =\frac{\eta(16\tau)^4}{\eta(8\tau)\eta(32\tau)}.
 \tag{3}\label{eq:theta}
\end{align}
Here $\Theta_{a,b,d}(\tau)=\sum_{x,y\in\Z}q^{ax^2+bxy+dy^2}$.
The completed function is entire, satisfies
$\Lambda(s)=\Lambda(1-s)$, and obeys
\[
 \Lambda(s)>0\quad(s\in\mathbb R),\qquad L(s)>0\quad(s>0).
 \tag{4}\label{eq:positive}
\]
Its only real zeros before completion are the simple trivial zeros of
$L(s)$ at $s=0,-1,-2,\ldots$.
\end{theorem}

For the arithmetic statement, let $p$ be any rational prime with
$p\equiv5\pmod8$. Let $Q_n$ be the degree-$p^n$ layer of the
cyclotomic $\Z_p$-extension of $\Q$, and set
$F_n=FQ_n$, $E_n=EQ_n$, and $E_\infty=\bigcup_n E_n$.
Put $A(k)=\Cl(k)\{p\}$ and $\Gamma=\Gal(E_\infty/E)$.
Since $[E:\Q]=8$, the composita are linearly disjoint from the
cyclotomic tower and
$\Gal(E_\infty/K)=\Delta\times\Gamma$.

Let $V/\Q_p$ be finite, with integer ring $\mathcal O$, and let
$\psi:\Delta\longrightarrow\mu_4\subset\Z_p^\times$ be faithful.
For any continuous $\kappa:\Gamma\longrightarrow\mathcal O^\times$, put
$D=(V/\mathcal O)(\psi\kappa)$, with action inflated through
$\Delta\times\Gamma$. The strict condition is
\[
 \Sel_{\rm str}(K,D)=
 \ker\left(H^1(K,D)\longrightarrow
 \prod_{v\mid p}H^1(K_v,D)\times
 \prod_{\substack{v\nmid p\\v\ {\rm finite}}}H^1(I_v,D)\right).
 \tag{5}\label{eq:strict}
\]

\begin{theorem}\label{thm:arithmetic}
One has $\Cl(F)=0$. For every rational prime $p\equiv5\pmod8$ and
every $n\ge0$,
\[
 A(F_n)=0,\qquad A(E_n)^\psi=0.
 \tag{6}\label{eq:class-zero}
\]
In particular, $\bigl(\varprojlim_n A(E_n)\bigr)^\psi=0$.
For every $V$, $\psi$, and continuous $\kappa$ as above,
\[
 \Sel_{\rm str}(K,(V/\mathcal O)(\psi\kappa))=0.
 \tag{7}\label{eq:selmer-zero}
\]
This includes all integral or $p$-adic principal weights and all finite
wild cyclotomic twists. Its normalized logarithmic size is zero under
every coefficient-field normalization.
\end{theorem}

The theorem concerns the two faithful characters of $\Delta$.
It does not identify the conjugation-odd part of the trivial or quadratic
character components. The real nonvanishing statement in
Theorem~\ref{thm:analytic} is separate from the class-group proof.
The detailed arithmetic proof first treats $p=5$ and then proves the
entire prime family in Corollary~\ref{cor:prime-family}.

\subsection{The cubic comparison and the scope of positivity}

There is a classical comparison in a different imaginary quadratic
field. Put $K_{23}=\Q(\sqrt{-23})$, let $H_{23}$ be its Hilbert class
field, and let $\vartheta$ be either nontrivial character of
$\Gal(H_{23}/K_{23})$. The theta identity below is stated explicitly
by \citet[p.~5 of the author version]{DukeLi2015}; the corresponding
Frobenius traces are also identified by
\citet[Theorem~3.1, p.~14]{deShalit2012}.

\begin{proposition}\label{prop:23}
The class group of $K_{23}$ is cyclic of order three, and $H_{23}$ is
the splitting field of $X^3-X-1$. If
\[
 L_{23}(s)=L_{K_{23}}(s,\vartheta),\qquad
 \Lambda_{23}(s)=23^{s/2}(2\pi)^{-s}\Gamma(s)L_{23}(s),
 \tag{32}\label{eq:23-completion}
\]
then $\Lambda_{23}$ is entire, satisfies
$\Lambda_{23}(s)=\Lambda_{23}(1-s)$, and is strictly positive for
every real $s$. In particular, $L_{23}(s)>0$ for $s>0$; its only
real zeros are precisely the simple trivial zeros at nonpositive integers.
\end{proposition}



For precision, a positive Mellin kernel together with a functional
equation about $1/2$ does not by itself force all zeros onto that line.
Here is an explicit entire counterexample to that implication. It is
an analytic example, not an arithmetic $L$-function.

\begin{proposition}\label{prop:kernel-limit}
There is a nonnegative smooth compactly supported Mellin kernel whose
entire transform $\mathcal M(s)$ is positive for every real $s$ and
satisfies $\mathcal M(s)=\mathcal M(1-s)$, but has zeros on both lines
$\Re s=1/4$ and $\Re s=3/4$.
\end{proposition}

\section*{Proof overview}

Because the Gaussian integers are principal, the Hecke character can be
written on a generator $a+bi$ by one residue modulo eight. Conjugate
generators cancel its imaginary values. Its remaining coefficients
separate into an odd-square theta function and an alternating-square
theta function. Poisson summation turns the latter into a sum of positive
Gaussians. The same calculation gives the exact transformation at the
two endpoints of the Mellin integral, proving convergence and positivity
for every real argument.

The arithmetic proof starts with an integral basis, a Minkowski bound,
and three polynomial factorizations. They prove the base class number
is one. Quadratic reciprocity and a quadratic-factor comparison show
that every prime $p\equiv5\pmod8$ is inert in the quartic field.
A single totally ramified prime identifies
class-group coinvariants in each finite cyclotomic layer with the base
class group. A nilpotent-operator argument forces their vanishing.
Descent of degree prime to $p$ and complex conjugation isolate the two faithful
quartic characters. Finally every Selmer cocycle has finite image, so
finite-level inflation and restriction prove the last assertion without
an infinite-layer control assumption.

\section{The explicit ring class character}

The norm on $\Z[i]$ is Euclidean. Hence $K$ has class number one.
The ring class group of $R=\Z+8\Z[i]$ is
\[
 (\Z[i]/8\Z[i])^\times\big/
 \bigl((\Z/8\Z)^\times\mu_4\bigr).
 \tag{8}\label{eq:ring-class}
\]
There are $32$ numerator elements and $8$ elements in the denominator,
so this quotient has order four. Its cyclic structure and character can
be seen without a class-group algorithm. For $a+bi$ of odd norm,
multiply by a unit so that $a$ is odd and $b$ is even, and define
\[
 t\equiv ba^{-1}\pmod8,\quad t\in\{0,2,4,6\},\qquad
 \chi((a+bi))=i^{t/2}.
 \tag{9}\label{eq:character}
\]
The two possible normalized generators differ by $-1$, which leaves
$t$ unchanged. Multiplication gives the rule
\[
 t\star u\equiv(t+u)(1-tu)^{-1}\pmod8.
 \tag{10}\label{eq:law}
\]
In fact $t\star u\equiv t+u\pmod8$: write $t=2r$, $u=2s$ and note
that $4rs(t+u)$ is divisible by eight. Thus
$t\mapsto i^{t/2}$ identifies the quotient in \eqref{eq:ring-class}
with $\mu_4$. Its kernel consists precisely of unit multiples of odd
rational residues.

The character is trivial on elements congruent to one modulo $(8)$.
It is nontrivial on $5+4i$, for which
\[
 \chi((5+4i))=-1,\qquad (5+4i)-1=4(1+i).
 \tag{11}\label{eq:conductor-witness}
\]
The ideal $(8)$ is the sixth power of the prime $(1+i)$, up to a unit,
and $4(1+i)$ has valuation five at that prime. Consequently the exact
conductor is $(8)$, not a proper divisor.
By the ring class field identification of
\citet[Lemma~3.8]{WilliamsLiu1994}, this is a faithful character of
$\Gal(E/K)$; the opposite Artin convention replaces it by its inverse.
All expressions for $L$ below are unchanged by that replacement.

One can also verify that field identification directly. The dihedral
permutation representation on the four roots is
$\boldsymbol1\oplus\chi_8\oplus\rho_E$.
Lemma~\ref{lem:base} below gives $D_F=-2048$, so the
conductor-discriminant formula gives conductor $256$ for $\rho_E$.
For either faithful character $\xi$ of $\Gal(E/K)$, induction gives
$256=|D_K|\N\mathfrak f(\xi)$; hence
$\mathfrak f(\xi)=(8)$. Every rational odd principal ideal has trivial
Artin symbol in $E/K$. At a rational prime split in $K$, the two symbols
are inverse under conjugation. At a prime inert in $K$, its rational
Frobenius is a dihedral reflection, whose square is one; that square is
the Frobenius over $K$. Thus the Artin kernel contains principal ideals
with generators congruent to odd rational integers modulo $(8)$.
The ring class correspondence embeds $E$ in the class field of $R$.
Both have degree four over $K$, so they are equal. This verifies the
classical identification using the displayed discriminant and the
usual Artin conductor and class-field correspondences.

For comparison, the reduced primitive positive forms of discriminant
$-256$ are exactly
\[
 (1,0,64),\quad(4,4,17),\quad(5,2,13),\quad(5,-2,13).
 \tag{12}\label{eq:forms}
\]
This is a finite enumeration: $1\le a\le\sqrt{256/3}$,
$|b|\le a\le d$, $b^2-4ad=-256$, with the usual nonnegative boundary
choice. The square of $(5,2,13)$ composes to $(25,12,4)$, and the
determinant-one substitution $(x,y)=(-v,u+2v)$ transforms this to
$(4,4,17)$. This also certifies a cyclic group of order four.

Every odd-norm Gaussian ideal has exactly two generators with odd real
part and even imaginary part. Therefore, for $\Re s>1$,
\[
 L(s)=\frac12\sum_{\substack{a\in\Z\ {\rm odd}\\b\in2\Z}}
       \frac{\chi((a+bi))}{(a^2+b^2)^s}.
 \tag{13}\label{eq:gaussian-series}
\]
Pairing $b$ with $-b$ cancels the purely imaginary character values.
If $b=4v$, the remaining value is $(-1)^v$. Thus
\[
 L(s)=\frac12\sum_{\substack{a\in\Z\ {\rm odd}\\v\in\Z}}
            \frac{(-1)^v}{(a^2+16v^2)^s}.
 \tag{14}\label{eq:separated-series}
\]
Absolute convergence in this half-plane justifies the regrouping.
This proves the separated theta expression in \eqref{eq:theta}.

For its binary-form expression, write
$4x^2+4xy+17y^2=(2x+y)^2+16y^2$.
The substitution $u=2x+y$ ranges over pairs with $u\equiv y\pmod2$.
The even-$y$ terms cancel the even-$u$ part of
$\Theta_{1,0,64}$, leaving exactly
$\sum_{u\ {\rm odd},v\in\Z}(-1)^vq^{u^2+16v^2}$.
Finally the classical Jacobi products
\citep[Equations~20.5.2 and 20.5.4]{DLMF} give
\[
 \frac12\sum_{u\ {\rm odd}}q^{u^2}
 =q\prod_{m\ge1}\frac{(1-q^{16m})^2}{1-q^{8m}},\qquad
 \sum_v(-1)^vq^{16v^2}
 =\prod_{m\ge1}\frac{(1-q^{16m})^2}{1-q^{32m}}.
 \tag{15}\label{eq:jacobi-products}
\]
Multiplication proves the eta quotient, including its leading factor $q$.

\section{Conductor, functional equation, and positivity}

The four-dimensional permutation representation on the roots of
$X^4-2$ decomposes as
\[
 \operatorname{Ind}_{\langle c\rangle}^{\Gal(E/\Q)}\boldsymbol1
   =\boldsymbol1\oplus\chi_8\oplus\rho,
 \qquad \rho=\operatorname{Ind}_\Delta^{\Gal(E/\Q)}\chi.
 \tag{16}\label{eq:permutation}
\]
Indeed $\chi_8$ is the character of the subfield $\Q(\sqrt2)$ and
$\rho$ is the unique irreducible representation of degree two of the
dihedral group. Matrices for $\rho(\delta)$ and $\rho(c)$ are a rotation
through a right angle and $\operatorname{diag}(1,-1)$, respectively.
Thus $\det\rho=\chi_{-4}$. Artin formalism proves
\eqref{eq:zeta-quotient}, including its ramified local factors.
The conductor induction formula gives
$\mathfrak f(\rho)=|D_K|\N(8)=4\cdot64=256$.
As an independent normalization check, the integral basis proved in the
next section gives $D_F=-2048$, and the conductor-discriminant formula
for \eqref{eq:permutation} gives $|D_F|=8\cdot256$.
These identifications also follow from Ishii's Artin construction
\citep[\S2, pp.~119--120, and \S3, Table~3]{Ishii1985}.

Here is a direct analytic proof, including the functional equation.
For $t>0$ put
\[
 A(t)=\frac12\sum_{u\ {\rm odd}}e^{-2\pi t u^2},\qquad
 B(t)=\sum_{v\in\Z}(-1)^v e^{-32\pi t v^2},
 \qquad f(it)=A(t)B(t).
 \tag{17}\label{eq:AB}
\]
Poisson summation for the Gaussian gives
\[
 B(t)=\frac1{\sqrt{32t}}
       \sum_{m\in\Z}e^{-\pi(m+1/2)^2/(32t)}>0.
 \tag{18}\label{eq:B-positive}
\]
Every summand in $A(t)$ is positive, so $f(it)>0$ for every $t>0$.
The two Gaussian transforms more precisely give, with $t'=1/(256t)$,
\[
 A(t)=\frac{B(t')}{4\sqrt{2t}},\qquad
 B(t)=\frac{A(t')}{2\sqrt{2t}},\qquad
 f(it)=\frac1{16t}f\left(\frac{i}{256t}\right).
 \tag{19}\label{eq:theta-transform}
\]
For example, apply
$\sum_n e^{-\pi a(n+1/2)^2}
=a^{-1/2}\sum_m(-1)^m e^{-\pi m^2/a}$
to $A(t)=\frac12\sum_n e^{-8\pi t(n+1/2)^2}$.
The ordinary Fourier transform
$\widehat{e^{-\pi a x^2}}(\xi)=a^{-1/2}e^{-\pi\xi^2/a}$
proves both displayed Poisson identities with these constants.

As $t\to\infty$, the separated series is $O(e^{-2\pi t})$.
Equation~\eqref{eq:theta-transform} therefore gives
$f(it)=O(t^{-1}e^{-\pi/(128t)})$ as $t\to0^+$.
Consequently
\[
 \Lambda(s)=16^s\int_0^\infty f(it)t^{s-1}\,dt
 \tag{20}\label{eq:Mellin}
\]
converges locally uniformly for every $s\in\C$ and defines an entire
function. For $\Re s>1$, termwise integration of the absolutely
convergent lattice series gives exactly \eqref{eq:completion}.
The substitution $t=1/(256u)$ in \eqref{eq:Mellin}, followed by
\eqref{eq:theta-transform}, proves
$\Lambda(s)=\Lambda(1-s)$.
For real $s$, the integrand and the factor $16^s$ are strictly positive.
Since $\Gamma(s)>0$ for $s>0$, this proves \eqref{eq:positive}.
Finally $1/\Gamma(s)$ has simple zeros exactly at the nonpositive
integers and no other zeros. Equation~\eqref{eq:completion}, with
$\Lambda(s)>0$ on the real axis, proves the last assertion of
Theorem~\ref{thm:analytic}.

\section{Class number one and the unique prime above five}

For the detailed arithmetic argument through the beginning of
Section~6, specialize the notation of Section~1 to $p=5$.
Corollary~\ref{cor:prime-family} then proves the assertion for every
prime in Theorem~\ref{thm:arithmetic}.

\begin{lemma}\label{lem:base}
The ring of integers of $F$ is $\Z[\alpha]$, its discriminant is $-2048$,
and its signature is $(2,1)$. Its class number is one. The prime five
is unramified and inert in $F$.
\end{lemma}
\begin{proof}
The polynomial $X^4-2$ is Eisenstein at two and has discriminant
$-2^{11}$. Thus only two could divide the index of $\Z[\alpha]$.
For Dedekind's index criterion its reduction is $X^4$ and
\[
 \frac{X^4-(X^4-2)}2=1.
 \tag{21}\label{eq:index}
\]
The required greatest common divisor with the repeated factor $X$ is
one. Hence two does not divide the index, proving the asserted basis
and discriminant. The two real roots and one nonreal conjugate pair
give the signature.

Minkowski's bound \citep[Theorem~4.3]{Milne2020} is
\[
 \frac{4!}{4^4}\frac4\pi\sqrt{2048}
   =\frac{12\sqrt2}{\pi}<6.
 \tag{22}\label{eq:Minkowski}
\]
For the strict inequality it suffices to use $\pi>3$ and $288<324$.
Only primes above $2,3,5$ could have norm below six. Their defining
polynomial factorizations are
\[
 \begin{array}{c|c}
 2&X^4\\
 3&(X^2+X+2)(X^2-X+2)\\
 5&X^4-2\quad\text{irreducible}.
 \end{array}
 \tag{23}\label{eq:factorizations}
\]
The two quadratic factors at three have nonsquare discriminant $2$.
At five, reduction modulo $X^4-2$ gives $X^{25}=4X$ and $X^{625}=X$.
Thus the polynomial divides $X^{625}-X$ and has no factor of degree
one or two, proving irreducibility. The only prime ideal with norm
below six is consequently the unique prime above two. It is
$(\alpha)$, of norm two, and is principal. Every ideal class has an
integral representative of norm below six, so every class is trivial.
The factorization at five also proves its inertness and, since
$5\nmid2048$, its absence of ramification.
\end{proof}

\section{A finite class-field argument in the tower}

\begin{lemma}\label{lem:coinvariants}
Let $L/k$ be cyclic of degree $p^r$ for an odd prime $p$. Suppose exactly
one finite prime of $k$ ramifies, and it is totally ramified. Then the
norm map induces an isomorphism
\[
 A_p(L)/(\sigma-1)A_p(L)\simeq A_p(k),
 \qquad A_p(k)=\Cl(k)\{p\}.
 \tag{24}\label{eq:coinvariants}
\]
If $A_p(k)=0$, then $A_p(L)=0$.
\end{lemma}
\begin{proof}
Let $H$ be the Hilbert $p$-class field of $L$.
It is Galois over $k$. Put $G=\Gal(H/k)$ and
$A=\Gal(H/L)=A_p(L)$. An inertia subgroup $I$ at the ramified prime
has order $p^r$, because $H/L$ is unramified, and maps isomorphically
onto $\Gal(L/k)$. Thus $G=A\rtimes I$. As $A$ and $I$ are abelian,
$[G,G]=(\sigma-1)A$. Killing inertia therefore gives the abelian
quotient $A/(\sigma-1)A$, whose extension of $k$ is unramified
everywhere. There is no infinite-place ramification of odd order.

Conversely the Hilbert $p$-class field $H_k$ of $k$ is disjoint from
$L$: every nontrivial subextension of $L/k$ is ramified at the specified
prime. Its compositum $H_kL/L$ is unramified and abelian, hence lies
in $H$. Therefore the preceding quotient gives all of $H_k$.
Restriction of Artin automorphisms is the ideal norm, proving
\eqref{eq:coinvariants}.

If the quotient is zero, put $T=\sigma-1$. Then $A=TA$. On $A/pA$,
the relation $\sigma^{p^r}=1$ gives $T^{p^r}=0$.
Surjectivity of $T$ and its nilpotence force $A/pA=0$.
Since $A$ is a finite $p$-group, this implies $A=0$.
\end{proof}

The cyclotomic extension $Q_n/\Q$ is unramified outside five and
totally ramified at five. Lemma~\ref{lem:base} says the completion of
$F$ at its unique five-adic prime is the unramified quartic extension
of $\Q_5$. Composing with a totally ramified extension preserves its
degree and total ramification. Thus $F_n/F$ is cyclic of degree $5^n$,
with exactly one ramified prime, totally ramified.
Lemma~\ref{lem:coinvariants} and $A(F)=0$ give $A(F_n)=0$ for every $n$.

Let $C_n=A(E_n)$. Extension and norm for $E_n/F_n$ satisfy
\[
 Nj=2,\qquad jN=1+c.
 \tag{25}\label{eq:prime-to-five}
\]
Since two is invertible on five-primary groups, extension identifies
$A(F_n)$ with $C_n^{c=1}$. Hence $C_n^{c=1}=0$.
The four characters of $\Delta$ take values in $\Z_5$, so their
idempotents are integral at five. Complex conjugation sends $C_n^\psi$
to $C_n^{\psi^{-1}}$, and the two faithful characters are distinct.
For $x\in C_n^\psi$, the element $x+cx$ is conjugation-fixed and hence
zero. Projecting it to its $\psi$-component gives $x=0$.
This proves \eqref{eq:class-zero}, for either faithful character, and
also proves the assertion about inverse limits.

\section{Every strict cyclotomic twist}

We give the finite-coefficient argument in detail. The prime five splits
in $K$, since $X^2+1=(X-2)(X+2)$ over $\F_5$.
At each of these two primes, $E/K$ is the unramified quartic extension:
its defining polynomial $X^4-2$ is irreducible modulo five.
Consequently the decomposition group of $E_n/K$ at either five-adic
prime is the whole product
\[
 G_n=\Delta\times\Gal(Q_n/\Q).
 \tag{26}\label{eq:decomp}
\]
Outside five, the only possible ramification of $E_n/K$ is at two,
and its inertia has order four. To see that it is the whole $\Delta$,
note that $E/K$ is abelian, unramified outside two, and $K$ has class
number one. Quotienting by two-adic inertia would otherwise give a
nontrivial everywhere unramified abelian extension of $K$.

Fix a uniformizer $\varpi$ of $\mathcal O$ and put $D_m=D[\varpi^m]$.
There is an $n$ such that $\kappa$ modulo $\varpi^m$ is trivial on
$\Gal(E_\infty/E_n)$, so $D_m$ is trivial over $E_n$.
The endomorphism $\psi(\delta)-1$ is a unit on $D_m$.
Averaging over $\Delta$ therefore gives
\[
 H^j(G_n,D_m)=0\quad(j\ge0).
 \tag{27}\label{eq:vanishing-cohom}
\]
For clarity, $D_m^\Delta=0$, positive cohomology of $\Delta$ vanishes
because its order is invertible, and the Hochschild--Serre sequence
for $\Delta\subset G_n$ proves \eqref{eq:vanishing-cohom}.
Inflation and restriction consequently identify
\[
 H^1(K,D_m)\simeq
 \operatorname{Hom}_{\rm cts}(G_{E_n},D_m)^{G_n}.
 \tag{28}\label{eq:finite-descent}
\]
The same statement holds at both five-adic primes by
\eqref{eq:decomp}. Away from five, the relative inertia group has
order one or four, prime to five, so restriction on inertia cohomology
is injective. Thus \eqref{eq:finite-descent} preserves exactly the
strict local conditions.

Let $A'(E_n)$ denote $A(E_n)$ modulo the subgroup generated by primes
above five. Finite global class field theory now identifies the strict
group for $D_m$ with
\[
 \operatorname{Hom}(A'(E_n),D_m)^{G_n}=0.
 \tag{29}\label{eq:finite-Selmer}
\]
Indeed, full local vanishing at five requires complete splitting, and
inertia vanishing elsewhere requires an unramified abelian extension.
Equivariance on $\Delta$ selects the $\psi$-component, which is zero
by \eqref{eq:class-zero}; the quotient $A'(E_n)$ has the same vanishing.

It remains to justify passage to $D$, including its local conditions.
A continuous cocycle from the compact group $G_K$ to the discrete
module $D$ has finite image, hence takes values in some $D_m$.
For any local group $J$, the exact sequence
$0\to D_m\to D\xrightarrow{\varpi^m}D\to0$ shows that the kernel of
\[
 H^1(J,D_m)\longrightarrow H^1(J,D)
 \tag{30}\label{eq:injection}
\]
is $D^J/\varpi^m D^J$. At either five-adic decomposition group,
$D^J=0$ by the faithful $\Delta$-action. At two-adic inertia it is
also zero. At every other inertia group $D^J=D$, which is divisible.
Thus \eqref{eq:injection} is injective for every local restriction
appearing in \eqref{eq:strict}. A strict class with values in $D$
therefore lifts to a strict class with values in $D_m$ and vanishes
by \eqref{eq:finite-Selmer}. This completes the proof of
Theorem~\ref{thm:arithmetic}.

\begin{corollary}\label{cor:prime-family}
Let $p\equiv5\pmod8$ be any rational prime. Let $Q_{p,n}$ be the
degree-$p^n$ layer of the cyclotomic $\Z_p$-extension of $\Q$, and set
$F_{p,n}=FQ_{p,n}$, $E_{p,n}=EQ_{p,n}$, and
$\Gamma_p=\Gal(EQ_{p,\infty}/E)$.
For either faithful character
$\psi:\Delta\to\mu_4\subset\Z_p^\times$, one has
\[
 \Cl(F_{p,n})\{p\}=0,\qquad
 \Cl(E_{p,n})\{p\}^\psi=0\qquad(n\ge0).
\]
For every finite extension $V/\Q_p$, with integer ring $\mathcal O$, and
every continuous $\kappa:\Gamma_p\to\mathcal O^\times$,
\[
 \Sel_{{\rm str},p}\bigl(K,(V/\mathcal O)(\psi\kappa)\bigr)=0.
\]
Here the strict condition is full decomposition-group vanishing at both
primes above $p$ and inertia vanishing at every other finite prime.
\end{corollary}
\begin{proof}
Quadratic reciprocity gives that $-1$ is a square modulo $p$, whereas
$2$ and $-2$ are nonsquares. The polynomial $X^4-2$ has no linear
factor. If it factored into two monic quadratics, write
\[
 X^4-2=(X^2+aX+b)(X^2-aX+d).
\]
Coefficient comparison gives
$b+d=a^2$, $a(d-b)=0$, and $bd=-2$.
If $a=0$, then $b^2=2$; if $a\ne0$, then $d=b$ and $b^2=-2$.
Both are impossible. Thus $p$ is inert in $F$ and, since $p>2$,
unramified. It splits in $K$ and has unramified degree four in $E/K$
at each of the two primes.

The base class number is one independently of $p$.
Lemma~\ref{lem:coinvariants} therefore gives
$\Cl(F_{p,n})\{p\}=0$. Because $p$ is odd, the extension-norm
identities \eqref{eq:prime-to-five}, with $p$-primary groups, again
identify the conjugation-fixed part with this zero group. Since
$\mu_4\subset\Z_p$ and four is invertible, the two faithful
components vanish separately. At the two $p$-adic places the full
finite decomposition group is $\Delta\times\Gal(Q_{p,n}/\Q)$;
at two the inertia is $\Delta$; elsewhere inertia acts trivially
on the coefficient module. Finally $\psi(\delta)-1$ is a $p$-adic
unit. Every cohomological vanishing and local injectivity in
\eqref{eq:vanishing-cohom}--\eqref{eq:injection} therefore holds with
$p$ replacing five, for every stated coefficient field and character.
This proves the strict-group assertion.
\end{proof}

\section{The unramified cubic character over
\texorpdfstring{$\mathbb Q(\sqrt{-23})$}{Q(sqrt(-23))}}


\begin{proof}[Proof of Proposition~\ref{prop:23}]
The reduced primitive forms of discriminant $-23$ are
$(1,1,6)$ and $(2,\pm1,3)$, so the class number is three.
The polynomial $X^3-X-1$ is irreducible by the rational-root test and
has squarefree discriminant $-23$. Its splitting field has group $S_3$
and quadratic subfield $K_{23}$. It is unramified outside $23$.
At $23$ the factorization is
\[
 X^3-X-1=(X-10)^2(X-3)\pmod{23}.
 \tag{33}\label{eq:23-factor}
\]
The tame inertia has order two and maps isomorphically to the inertia
of $K_{23}/\Q$. The cubic extension over $K_{23}$ is therefore
unramified also at $23$. Its degree three identifies it with $H_{23}$.

For $q=e^{2\pi i\tau}$, the classical normalized theta identity is
\[
 g(\tau)=\sum_{n\ge1}b_nq^n
 =\frac{\Theta_{1,1,6}(\tau)-\Theta_{2,1,3}(\tau)}2
 =\eta(\tau)\eta(23\tau),\qquad
 L_{23}(s)=\sum_{n\ge1}\frac{b_n}{n^s}.
 \tag{34}\label{eq:23-theta}
\]
The two nonprincipal ideal classes are inverses, their theta series
coincide, and $\vartheta+\vartheta^{-1}=-1$, which fixes the factor
$1/2$ in the binary-form expression. The eta equality is the published
identity just cited. Its leading coefficient is one. The Hecke character
is unramified, so its induced conductor is $|D_{K_{23}}|=23$.

For $t>0$, the defining product
\[
 \eta(it)=e^{-\pi t/12}\prod_{n\ge1}(1-e^{-2\pi nt})
 \tag{35}\label{eq:eta-positive}
\]
is positive: every factor is positive and the sum of the omitted
quantities converges. Thus $g(it)>0$. The eta transformation
$\eta(-1/\tau)=(-i\tau)^{1/2}\eta(\tau)$
\citep[Equations~23.18.5--23.18.7]{DLMF} gives
\[
 g(it)=\frac1{\sqrt{23}\,t}
             g\left(\frac{i}{23t}\right).
 \tag{36}\label{eq:23-transform}
\]
There is no phase ambiguity on the imaginary axis, where the square
roots are positive. At infinity $g(it)=O(e^{-2\pi t})$; at zero
\eqref{eq:23-transform} gives
$g(it)=O(t^{-1}e^{-2\pi/(23t)})$. Consequently
\[
 \Lambda_{23}(s)=23^{s/2}\int_0^\infty g(it)t^{s-1}\,dt
 \tag{37}\label{eq:23-Mellin}
\]
is entire and positive for real $s$. For $\Re s>1$ it equals
\eqref{eq:23-completion} by termwise integration. Substitution
$t=1/(23u)$ proves the functional equation, and division by
$\Gamma(s)$ gives the assertions about the uncompleted function.
\end{proof}

No cyclotomic Selmer claim for $K_{23}$ is required for this proposition.
The two theta products establish two explicit real nonvanishing results.

\section{What positivity and a functional equation do not imply}


\begin{proof}[Proof of Proposition~\ref{prop:kernel-limit}]
Let
\[
 \phi(u)=
 \begin{cases}
  \exp\bigl(-1/(1-16u^2)\bigr),& |u|<1/4,\\
  0,& |u|\ge1/4,
 \end{cases}
 \qquad a=\cosh(1/4)>1,
\]
and set $m(u)=a\phi(u)+\tfrac12\phi(u-1)+\tfrac12\phi(u+1)$.
This is even, nonnegative, smooth, compactly supported, and nonzero.
With $\Phi(z)=\int_{\mathbb R}\phi(u)e^{zu}\,du$, direct substitution
gives
\[
 \mathcal M(s)=\int_{\mathbb R}m(u)e^{(s-1/2)u}\,du
  =\Phi(s-1/2)\bigl(a+\cosh(s-1/2)\bigr).
 \tag{38}\label{eq:kernel-counterexample}
\]
Compact support proves entire continuation. Evenness proves the
functional equation, and positivity of the integrand proves positivity
on the real axis. Since
\[
 \cosh\bigl(\pm1/4+(2k+1)\pi i\bigr)=-\cosh(1/4),
\]
the points $s=1/4+(2k+1)\pi i$ and $s=3/4+(2k+1)\pi i$ are zeros.
Finally the substitution $x=e^u$ writes the transform as
$\int_0^\infty m(\log x)x^{-1/2}x^{s-1}\,dx$, exhibiting the stated
positive Mellin kernel on $(0,\infty)$.
\end{proof}

There is no Euler-product assertion for \eqref{eq:kernel-counterexample}.
The example isolates the logical scope of the positivity argument;
it says nothing against RH for the arithmetic functions above.

\section{Finite checks and scope}

Every arithmetic input to the proof is displayed in
\eqref{eq:character}--\eqref{eq:forms} and
\eqref{eq:index}--\eqref{eq:factorizations}.
For example, a direct calculation on the $32$ units modulo eight
checks all $1024$ products in \eqref{eq:character}. Each of its four
values occurs eight times. No conjectural class-group computation is
used at any level.

An exact integer calculation also compared four independently evaluated
coefficient constructions through $q^{4096}$: the two binary theta
series, the separated signed theta sum, the eta product, and the Euler
coefficients recovered from $X^4-2$ modulo odd primes. This finite
comparison checks the implementations and normalizations; the infinite
identities were proved above. The first nonzero terms are
\[
 f(\tau)=q+q^9-2q^{17}-q^{25}-2q^{41}+q^{49}
          +2q^{73}+q^{81}+2q^{89}-2q^{97}
          +2q^{113}+q^{121}+O(q^{130}).
 \tag{31}\label{eq:first-coefficients}
\]
For the Euler construction, at an odd prime $p$ let $r_p$ be the number
of roots of $X^4-2$ in $\F_p$. Then
$a_p=r_p-1-\chi_8(p)$ and
$a_{p^j}=a_pa_{p^{j-1}}-\chi_{-4}(p)a_{p^{j-2}}$.
At two the Euler factor is one. Multiplicativity determines the
coefficients compared above.
For the $-23$ comparison, the three reduced forms, the cubic
discriminant, the factorization \eqref{eq:23-factor}, and equality of
the binary-theta and eta-product coefficients through $q^{4096}$ were
also checked by exact arithmetic.

The analytic results prove real nonvanishing of the two specified Hecke
$L$-functions. They provide no zero-free half-plane beyond the real-axis
assertion and no RH statement. The arithmetic theorem proves the full
vanishing of the stated strict groups, not merely a bound on their
exponents or lengths. It requires a faithful quartic character and
the strict local conditions in \eqref{eq:strict}.

\section*{Method and availability}
The mathematics, text, and original computations were produced by GPT-6
Astra in ChatGPT Pro mode. The classical theta and ring-class identities
are credited to Ishii, Williams and Liu, Duke and Li, and de Shalit;
the class-number and theta-product inputs are cited separately.
The character normalization, positivity argument, finite class-field
control, and strict cyclotomic descent are proved here. Exact arithmetic
checks concern the displayed finite witnesses. The accompanying records
state the scope of targeted source and proof checks by the same model.
No human review or formal verification is claimed.
The original text and programs are dedicated to the public domain under
CC0 1.0; cited works retain their own terms. Sources and verification
records are at \url{\workbenchurl}, with version records at
\url{\recordurl}.

\begin{thebibliography}{9}
\bibitem[de Shalit(2012)]{deShalit2012}
de Shalit, E. (2012, August).
\emph{Squares: An invitation to modern number theory}.
Expanded notes for the Asian Science Camp.
\url{https://ma.huji.ac.il/~deshalit/new_site/files/squares.pdf}.

\bibitem[Duke and Li(2015)]{DukeLi2015}
Duke, W., \& Li, Y. (2015). Harmonic Maass forms of weight 1.
\emph{Duke Mathematical Journal}, \emph{164}(1), 39--113.
Author version dated March 2, 2014, p.~5.
\url{https://www.math.ucla.edu/~wdduke/preprints/weight%20one.pdf}.

\bibitem[Fukuda(1994)]{Fukuda1994}
Fukuda, T. (1994). Remarks on $\mathbb Z_p$-extensions of number fields.
\emph{Proceedings of the Japan Academy, Series A, Mathematical Sciences},
\emph{70}(8), 264--266.
\url{https://doi.org/10.3792/pjaa.70.264}.

\bibitem[Ishii(1985)]{Ishii1985}
Ishii, N. (1985). Cusp forms of weight one, quartic reciprocity and
elliptic curves. \emph{Nagoya Mathematical Journal}, \emph{98}, 117--137.
\url{https://doi.org/10.1017/S0027763000021413}.

\bibitem[Milne(2020)]{Milne2020}
Milne, J. S. (2020). \emph{Algebraic number theory} (Version 3.08).
\url{https://www.jmilne.org/math/CourseNotes/ANT.pdf}.

\bibitem[National Institute of Standards and Technology(2026)]{DLMF}
National Institute of Standards and Technology. (2026).
\emph{Digital Library of Mathematical Functions} (Version 1.2.8).
Sections 20.5 and 23.18.
\url{https://dlmf.nist.gov/20.5}; \url{https://dlmf.nist.gov/23.18}.

\bibitem[Williams and Liu(1994)]{WilliamsLiu1994}
Williams, K. S., \& Liu, D. (1994).
Representation of primes by the principal form of negative discriminant
$\Delta$ when $h(\Delta)$ is 4.
\emph{Tamkang Journal of Mathematics}, \emph{25}(4), 321--334.
\url{https://journals.math.tku.edu.tw/index.php/TKJM/article/view/4461}.
\end{thebibliography}
\end{document}
